\section{Objetivo y contexto de la evaluación}
\label{sec:contexto}

Este informe presenta la evaluación de línea base de
\NombreProyecto, correspondiente al cuarto entregable del
Desafío de Investigación de \CodigoCurso\ \citep{enunciado2026}.

El problema elegido es resolver Sudokus de tamaño $N \times N$: completar
una grilla de modo que cada fila, columna y bloque contenga los números
de 1 a $N$ exactamente una vez, sin modificar las pistas originales. Es
un problema de razonamiento con restricciones, de varios pasos, cuya
respuesta se puede verificar con código sin ambigüedad. No requiere
conocimiento externo ni herramientas, y su dificultad se puede controlar
con el tamaño del tablero y la cantidad de casillas vacías. El modelo
objetivo es Qwen3.5-2B, que resuelve con dificultad esta tarea.

La evaluación busca establecer, antes de cualquier entrenamiento, cuánto
resuelve el modelo objetivo en cada tamaño y nivel de dificultad, cómo se
compara con un modelo mayor de la misma familia y si las tareas caen en
una zona aprendible: niveles en que el modelo acierta en algunos intentos
y falla en otros, condición necesaria para que el aprendizaje reforzado
con recompensas verificables tenga señal. No se entrenaron pesos en esta
etapa.

El harness define el formato del prompt, el parser de la respuesta y el
verificador, expuesto mediante la interfaz común
\texttt{reward(prompt, completion) -> float}. Con el mismo verificador se
generaron también los datos derivados del profesor Qwen3.5-9B: 161
ejemplos SFT y 600 pares de preferencia DPO a partir de la partición de
entrenamiento. El código, el dataset y los datos derivados se describen
en sus entregables respectivos; aquí se resume solo lo necesario para
interpretar la línea base.

\subsection{Tareas, particiones y dificultad}
Las tareas se generan con código, adaptando el generador de Sudoku de
Reasoning Gym \citep{stojanovski2025reasoninggym}: primero se construye
una solución completa aleatoria y luego se borran casillas una a una,
deshaciendo cada borrado que permitiría una segunda solución. Así, cada
puzzle tiene solución única y su respuesta es verificable por
construcción. Aun así, el verificador no compara con la solución de
referencia: acepta cualquier tablero que cumpla las reglas y conserve
las pistas.

El dataset contiene tableros de 4x4 (bloques de 2x2), 6x6 (bloques de
2x3) y 9x9 (bloques de 3x3), con tres niveles definidos por la cantidad
de casillas vacías (tabla~\ref{tab:niveles}). Cada combinación de tamaño
y nivel tiene 50 tareas de entrenamiento, 10 de validación y 10 de
prueba: en total, 450, 90 y 90 tareas. La línea base de este informe usa
las 90 tareas de validación.

\begin{table}[htbp]
  \centering\small
  \caption{Casillas vacías por tamaño de tablero y nivel. Los rangos se
  fijaron antes de la evaluación y no están calibrados con el modelo
  objetivo.}
  \label{tab:niveles}
  \begin{tabular}{@{}lccc@{}}
    \toprule
    Tablero & \emph{easy} & \emph{medium} & \emph{hard} \\
    \midrule
    4x4 (16 casillas) & 4--6   & 7--9   & 10--12 \\
    6x6 (36 casillas) & 8--12  & 13--18 & 19--24 \\
    9x9 (81 casillas) & 10--20 & 21--35 & 36--50 \\
    \bottomrule
  \end{tabular}
\end{table}

Los niveles miden la cantidad de casillas por completar, no la técnica
necesaria para resolverlas: 89 de las 90 tareas de validación se pueden
resolver completando repetidamente una casilla con un único candidato.
Los nombres \emph{easy}, \emph{medium} y \emph{hard} son, por lo tanto,
relativos a cada tamaño y no comparables entre tamaños.

Cada tarea se genera con su propia semilla, de modo que el dataset es
reproducible. Las particiones se generan en el orden prueba, validación
y entrenamiento, y ningún puzzle se repite entre ellas; aumentar el
número de tareas de entrenamiento no altera las otras dos. Una revisión
posterior de simetrías, limitada a los tableros de 4x4, comparó las
particiones de entrenamiento y validación bajo renombramiento de
números, intercambio de filas o columnas dentro de una banda, intercambio
de bandas y transposición. Aunque no hay puzzles idénticos, 3 de las 10
tareas de validación de 4x4 \emph{hard} resultaron equivalentes a tareas
de entrenamiento. Esto no afecta la línea base, que no involucra
entrenamiento, pero limitará las conclusiones sobre generalización en
el proyecto posterior; los tableros de 6x6 y 9x9 no se revisaron.

La partición de prueba no se usó en ninguna etapa: ni para ajustar las
tareas o el protocolo, ni para generar datos derivados. Las decisiones de
calibración se toman exclusivamente con la partición de validación, y la
de prueba queda reservada para la evaluación final.
